\documentclass[11pt, a4paper]{exam}
\usepackage[a4paper,top=2.3cm,bottom=2.8cm,left=2cm,right=2cm,headheight=15pt,headsep=0.6cm,footskip=1.2cm]{geometry}
\usepackage[T1]{fontenc}
\usepackage[utf8]{inputenc}
\usepackage{lmodern}
\usepackage{textcomp}
\usepackage[french]{babel}
\usepackage{xcolor}
\usepackage{listings}
\usepackage{graphicx}
\usepackage{array}
\usepackage{enumitem}
\usepackage{amsmath}
\usepackage{tikz}
\usetikzlibrary{positioning,shapes.multipart,arrows.meta}
\usepackage[most]{tcolorbox}

\definecolor{codegreen}{rgb}{0,0.6,0}
\definecolor{codegray}{rgb}{0.5,0.5,0.5}
\definecolor{codepurple}{rgb}{0.58,0,0.82}
\definecolor{backcolour}{rgb}{0.95,0.95,0.92}

\lstdefinestyle{mystyle}{
    backgroundcolor=\color{backcolour},
    commentstyle=\color{codegreen},
    keywordstyle=\color{magenta},
    numberstyle=\tiny\color{codegray},
    stringstyle=\color{codepurple},
    basicstyle=\ttfamily\footnotesize,
    breakatwhitespace=false,
    breaklines=true,
    postbreak=\mbox{\textcolor{codegray}{$\hookrightarrow$}\space},
    captionpos=b,
    keepspaces=true,
    columns=fullflexible,
    numbers=left,
    numbersep=5pt,
    showspaces=false,
    showstringspaces=false,
    showtabs=false,
    upquote=true,
    tabsize=4,
    xleftmargin=1.5em,
    framexleftmargin=0.3em,
    aboveskip=0.6em,
    belowskip=0.6em,
    extendedchars=true,
    literate={é}{{\'e}}1 {è}{{\`e}}1 {ê}{{\^e}}1 {ë}{{\"e}}1 {à}{{\`a}}1 {â}{{\^a}}1
             {ç}{{\c c}}1 {ù}{{\`u}}1 {û}{{\^u}}1 {ô}{{\^o}}1 {î}{{\^i}}1 {ï}{{\"i}}1
             {É}{{\'E}}1 {È}{{\`E}}1 {À}{{\`A}}1 {Ç}{{\c C}}1 {œ}{{\oe}}1
             {«}{{\guillemotleft}}1 {»}{{\guillemotright}}1 {→}{{$\rightarrow$}}1
}
\lstset{style=mystyle, language=Python}
\lstdefinestyle{sqlstyle}{style=mystyle, language=SQL, morekeywords={LIMIT,REFERENCES,IN,TEXT,FLOAT,DATE,VARCHAR,CHAR,REPLACE}, numbers=none, xleftmargin=0.3em, frame=lines}

% Code en ligne
\newcommand{\code}[1]{\texttt{#1}}

% Titre d'une question : numéro et intitulé, barème à droite
\newcommand{\Qt}[2]{\subsection*{#1\hfill{\normalfont\small #2}}}
% Rappel de la consigne en une ligne
\newcommand{\rappel}[1]{\noindent\textit{Ce qui est demandé :} #1\par\smallskip}

% La réponse attendue, mise en évidence
\newtcolorbox{reponse}[1][Réponse]{enhanced, breakable, colback=black!5, colframe=black, boxrule=0.9pt, arc=1mm,
  title={#1}, fonttitle=\bfseries\small, coltitle=black, colbacktitle=black!15, left=2mm, right=2mm, top=1mm, bottom=1mm,
  before skip=0.5em, after skip=0.7em}
% Erreur fréquente
\newtcolorbox{erreur}{enhanced, breakable, colback=white, colframe=black, boxrule=0.5pt, leftrule=3pt, arc=0mm,
  title={Erreur fréquente}, fonttitle=\bfseries\small, coltitle=black, colbacktitle=white, left=2mm, right=2mm, top=1mm, bottom=1mm,
  before skip=0.5em, after skip=0.7em}
% Encadré « À retenir »
\newtcolorbox{retenir}{enhanced, breakable, colback=white, colframe=black, boxrule=0.5pt, arc=1mm, boxsep=1mm,
  borderline={0.4pt}{1.2pt}{black, dashed},
  title={À retenir}, fonttitle=\bfseries\small, coltitle=black, colbacktitle=white, left=2mm, right=2mm, top=1mm, bottom=1mm,
  before skip=0.5em, after skip=0.7em}

\newcommand{\methode}{\noindent\textbf{Méthode.}\ }
\newcommand{\explication}{\noindent\textbf{Explication.}\ }

\setlist{itemsep=0.15em, topsep=0.3em}
\renewcommand{\arraystretch}{1.25}
\setlength{\parindent}{0pt}
\setlength{\parskip}{0.35em}

% Piles dessinées verticalement (sommet en haut) ; \suppr barre un élément supprimé
\newcommand{\suppr}[1]{\ooalign{#1\cr\hidewidth\rule[0.55ex]{1em}{1pt}\hidewidth\cr}}
\newcommand{\pilev}[2]{\begin{tikzpicture}[baseline=(titre.base)]
  \node[font=\scriptsize, align=center, text width=22mm] (titre) at (0,0) {#1};
  \foreach \v [count=\k] in {#2} {\node[draw, minimum width=8mm, minimum height=5mm, inner sep=0pt, font=\small] at (0,-0.25-\k*0.5) {\v};}
\end{tikzpicture}}
\setlength{\emergencystretch}{3em}
% Remarque sur l'énoncé (erreur ou ambiguïté du sujet)
\newtcolorbox{remarque}{enhanced, breakable, colback=white, colframe=black, boxrule=0.5pt, arc=0mm,
  borderline={0.4pt}{1.2pt}{black, dotted},
  title={Remarque sur l'énoncé}, fonttitle=\bfseries\small, coltitle=black, colbacktitle=black!8, left=2mm, right=2mm, top=1mm, bottom=1mm,
  before skip=0.5em, after skip=0.7em}
\usepackage{pgfplots}
\pgfplotsset{compat=1.17}
\newcommand{\sql}[1]{\textsc{#1}}
\lstdefinestyle{sqlc}{style=sqlstyle}
\newcommand{\resultat}[1]{\par\smallskip\noindent\textit{Résultat sur le jeu d'essai :}\par\nopagebreak\begin{center}\footnotesize #1\end{center}}

\title{Corrigé — CC2 Term NSI - SQL, Tri et complexité}
\author{Babash - NSI Tani Malandi et Johanne Bakalara}
\date{Contrôle de janvier 2024 (1 h)}

\begin{document}
\pagestyle{headandfoot}
\firstpageheader{}{}{}
\runningheader{\small Corrigé}{}{\small SQL, tri et complexité -- janvier 2024}
\runningheadrule
\firstpagefooter{}{\footnotesize Babash - NSI Tani Malandi et Johanne Bakalara \textperiodcentered{} CC BY-SA 4.0 \textperiodcentered{} page \thepage}{}
\runningfooter{}{\footnotesize Babash - NSI Tani Malandi et Johanne Bakalara \textperiodcentered{} CC BY-SA 4.0 \textperiodcentered{} page \thepage}{}

\maketitle
\thispagestyle{headandfoot}

\begin{tcolorbox}[colback=white, colframe=black, boxrule=0.5pt, arc=1mm, left=2mm, right=2mm]
\small\textbf{Comment lire ce corrigé.} Pour chaque question : ce qui est demandé, puis la réponse attendue (encadrée), puis l'explication pas à pas. Essaie d'abord seul·e, puis compare ta démarche, pas seulement ton résultat.
Les requêtes ont été exécutées avec SQLite et le code Python avec Python 3. Le sujet ne donne pas le contenu des tables : pour montrer un résultat, on utilise un petit \textbf{jeu d'essai} inventé.
\end{tcolorbox}

\begin{remarque}
Le barème du sujet donne 12 points à la partie SQL, sans détail par sous-question, puis 2 + 3 + 2 + 3 points : le total fait \textbf{22 points}. Ce corrigé reprend ces points ; c'est à l'enseignant·e de ramener la note sur 20 s'il le faut.
\end{remarque}

\section*{Partie 1 -- SQL (question 1, 12 points)}

\begin{retenir}
\begin{itemize}
\item Abonne(\underline{Abonne\_ID}, Nom, Prenom, DateNaissance, Adresse, Ville, CodePostal) ; Livre(\underline{Livre\_ID}, Titre, Auteur, Genre, ISBN) ; Emprunt(\#Abonne\_ID, \#Livre\_ID, DateEmprunt, DateRetourPrevu, DateRetourReel).
\item Les dates s'écrivent \code{'AAAA-MM-JJ'} : comparer ou trier ces textes revient à comparer les dates, car l'année vient en premier.
\end{itemize}
\end{retenir}

\textbf{Jeu d'essai inventé} (seules les colonnes utiles sont montrées ; les personnes sont fictives) :
\begin{center}\footnotesize
\begin{tabular}{|c|l|l|c|l|c|}\hline
\multicolumn{6}{|c|}{\textbf{Abonne}}\\\hline
Abonne\_ID & Nom & Prenom & DateNaissance & Ville & CodePostal \\\hline
801 & Abdou & Nassim & 2007-03-14 & Chirongui & 97620 \\
802 & Said & Fatima & 1958-11-02 & Chirongui & 97620 \\
803 & Ahmed & Inès & 1990-06-30 & Mamoudzou & 97600 \\
804 & Leroy & Paul & 1975-01-20 & Toulouse & 31000 \\
805 & Combo & Ali & 1949-08-09 & Dembéni & 97660 \\\hline
\end{tabular}\\[0.5em]
\begin{tabular}{|c|l|}\hline
\multicolumn{2}{|c|}{\textbf{Livre}}\\\hline
Livre\_ID & Titre \\\hline
4252 & Papillon \\
4253 & Le Lagon bleu \\\hline
\end{tabular}\quad
\begin{tabular}{|c|c|c|c|c|}\hline
\multicolumn{5}{|c|}{\textbf{Emprunt}}\\\hline
Abonne\_ID & Livre\_ID & DateEmprunt & DateRetourPrevu & DateRetourReel \\\hline
801 & 4252 & 2020-10-01 & 2020-10-22 & 2020-10-20 \\
804 & 4252 & 2020-10-25 & 2020-11-15 & NULL \\
802 & 4253 & 2020-10-05 & 2020-10-26 & 2020-10-26 \\\hline
\end{tabular}
\end{center}

\subsection*{Sélection}

\textbf{a)} \rappel{les noms et prénoms des abonné·es qui habitent à Chirongui.}
\begin{reponse}
\begin{lstlisting}[style=sqlc]
SELECT Nom, Prenom FROM Abonne WHERE Ville = 'Chirongui';
\end{lstlisting}
\end{reponse}
\resultat{\begin{tabular}{|l|l|}\hline Nom & Prenom \\\hline Abdou & Nassim \\ Said & Fatima \\\hline\end{tabular}}
\methode Quelles colonnes ? \code{Nom, Prenom}. Quelle table ? \code{Abonne}. Quelles lignes ? celles où \code{Ville} vaut \code{'Chirongui'}, entre guillemets car c'est un texte.

\textbf{b)} \rappel{le nom des personnes dont le code postal commence par 976.}
\begin{reponse}
\begin{lstlisting}[style=sqlc]
SELECT Nom FROM Abonne WHERE CodePostal LIKE '976%';
\end{lstlisting}
\end{reponse}
\resultat{\begin{tabular}{|l|}\hline Nom \\\hline Abdou \\ Said \\ Ahmed \\ Combo \\\hline\end{tabular}}
\explication \sql{LIKE} compare à un motif : \code{\%} remplace n'importe quelle suite de caractères. \code{'976\%'} veut dire « 976, puis n'importe quoi » : le code postal \textbf{commence} par 976 (les codes postaux de Mayotte).

\begin{erreur}
Écrire \code{'\%976\%'} : on cherche alors les codes qui \textbf{contiennent} 976, n'importe où. Ou écrire \code{CodePostal = '976\%'} : avec \code{=}, le \code{\%} n'est pas un joker, on cherche le texte exact « 976\% ».
\end{erreur}

\textbf{c)} \rappel{le nom des personnes qui ont emprunté le livre Papillon.}
\begin{reponse}
\begin{lstlisting}[style=sqlc]
SELECT Abonne.Nom
FROM Abonne
JOIN Emprunt ON Abonne.Abonne_ID = Emprunt.Abonne_ID
JOIN Livre ON Emprunt.Livre_ID = Livre.Livre_ID
WHERE Livre.Titre = 'Papillon';
\end{lstlisting}
\end{reponse}
\resultat{\begin{tabular}{|l|}\hline Nom \\\hline Abdou \\ Leroy \\\hline\end{tabular}}
\methode
\begin{enumerate}
\item Le nom est dans \code{Abonne}, le titre dans \code{Livre} ; le lien entre les deux est la table \code{Emprunt}.
\item Il faut donc \textbf{deux jointures} : \code{Abonne} avec \code{Emprunt} (par \code{Abonne\_ID}), puis \code{Emprunt} avec \code{Livre} (par \code{Livre\_ID}).
\item On garde les lignes dont le titre est \code{'Papillon'}.
\end{enumerate}
Si une même personne a emprunté le livre deux fois, son nom apparaît deux fois ; \code{SELECT DISTINCT Abonne.Nom} supprime les doublons.

\textbf{d)} \rappel{tous les attributs de l'abonné·e le plus âgé.}
\begin{reponse}
\begin{lstlisting}[style=sqlc]
SELECT * FROM Abonne ORDER BY DateNaissance ASC LIMIT 1;
\end{lstlisting}
\end{reponse}
\resultat{\begin{tabular}{|c|l|l|c|l|l|c|}\hline Abonne\_ID & Nom & Prenom & DateNaissance & Adresse & Ville & CodePostal \\\hline 805 & Combo & Ali & 1949-08-09 & 1 place de la Mairie & Dembéni & 97660 \\\hline\end{tabular}}
\explication La personne la plus âgée est celle qui est née \textbf{le plus tôt}, donc celle qui a la plus \textbf{petite} date de naissance. On trie par date croissante (\sql{ASC}) et on garde la première ligne (\sql{LIMIT} 1). \code{*} donne tous les attributs.

\textbf{Variante} avec une sous-requête, qui donne toutes les personnes en cas d'égalité :
\begin{lstlisting}[style=sqlc]
SELECT * FROM Abonne
WHERE DateNaissance = (SELECT MIN(DateNaissance) FROM Abonne);
\end{lstlisting}

\begin{erreur}
Trier par ordre décroissant (\sql{DESC}) ou utiliser \code{MAX(DateNaissance)} : on obtient la personne la plus \textbf{jeune}, née le plus tard.
\end{erreur}

\subsection*{Modification}

\textbf{e)} \rappel{ajouter l'emprunt du livre 4252 par l'abonné·e 803, le 2 novembre 2020, pour trois semaines.}
\begin{reponse}
\begin{lstlisting}[style=sqlc]
INSERT INTO Emprunt VALUES (803, 4252, '2020-11-02', '2020-11-23', NULL);
\end{lstlisting}
\end{reponse}
\methode
\begin{enumerate}
\item On \textbf{ajoute} une ligne : \sql{INSERT INTO}.
\item L'ordre des colonnes d'\code{Emprunt} est : \code{Abonne\_ID}, \code{Livre\_ID}, \code{DateEmprunt}, \code{Date\-Retour\-Prevu}, \code{Date\-Retour\-Reel}.
\item Trois semaines font 21 jours : 2 + 21 = 23, donc le retour est prévu le 2020-11-23 (novembre a 30 jours, on reste dans le mois).
\item Le livre n'est pas encore rendu : la date de retour réelle est inconnue, on met \sql{NULL} (« pas de valeur »).
\end{enumerate}
\textbf{Variante} en nommant les colonnes (la colonne non citée vaut \sql{NULL}) :
\begin{lstlisting}[style=sqlc]
INSERT INTO Emprunt (Abonne_ID, Livre_ID, DateEmprunt, DateRetourPrevu)
VALUES (803, 4252, '2020-11-02', '2020-11-23');
\end{lstlisting}

\begin{erreur}
Écrire les identifiants dans le désordre (\code{4252, 803, \ldots}) : la requête est acceptée, mais elle enregistre l'emprunt du livre 803 par l'abonné·e 4252. Autre erreur : écrire la date \code{'02/11/2020'} au lieu du format \code{'2020-11-02'} demandé.
\end{erreur}

\begin{remarque}
La table \code{Emprunt} n'a pas de clé primaire dans le schéma. Rien n'empêche donc d'enregistrer deux fois exactement la même ligne : nous l'avons vérifié, SQLite accepte la requête deux fois de suite. On pourrait ajouter un identifiant d'emprunt, ou prendre comme clé primaire le triplet (Abonne\_ID, Livre\_ID, DateEmprunt).
\end{remarque}

\subsection*{Ajout d'une table}

\textbf{f)} \rappel{créer la table \code{Emplacement} (id, zone, etagere).}
\begin{reponse}
\begin{lstlisting}[style=sqlc]
CREATE TABLE Emplacement (
    id INT PRIMARY KEY,
    zone VARCHAR(50),
    etagere INT,
    FOREIGN KEY (id) REFERENCES Livre(Livre_ID)
);
\end{lstlisting}
\end{reponse}
\explication
\begin{itemize}
\item Chaque attribut est suivi de son type : \code{INT} pour un entier, \code{VARCHAR(50)} pour un texte de 50 caractères au plus.
\item \code{id} est « le numéro du livre » : chaque livre a un seul emplacement, \code{id} peut donc être la clé primaire.
\item C'est aussi une clé étrangère vers \code{Livre(Livre\_ID)} : on ne peut pas ranger un livre qui n'existe pas. Cette ligne est un plus ; la consigne ne l'exige pas explicitement.
\end{itemize}

\section*{Partie 2 -- Algorithme de tri}

\begin{lstlisting}
def tri_surprise(tab):
    if len(tab) < 2:
        return tab
    for i in range(1, len(tab)):
        e = tab[i]
        j = i - 1
        while j >= 0 and e < tab[j] :
            tab[j+1] = tab[j]
            j = j - 1
        tab[j+1] = e
    return tab
\end{lstlisting}

\Qt{Question 2 -- Reconnaître l'algorithme et sa complexité}{2 points}

\textbf{a)} \rappel{quel est cet algorithme de tri ?}
\begin{reponse}
\textbf{Tri par insertion.}
\end{reponse}
\explication Pour chaque \code{i}, on prend l'élément \code{e = tab[i]} et on décale vers la droite les éléments déjà triés qui sont plus grands que lui, puis on \textbf{insère} \code{e} à la place libérée. Comme quand on range une carte dans une main déjà triée. Le tri par sélection, lui, cherche le minimum de la partie non triée et l'échange ; le tri fusion coupe le tableau en deux et s'appelle lui-même.

\textbf{b)} \rappel{la complexité de cet algorithme.}
\begin{reponse}
\textbf{Quadratique}, en $O(n^2)$.
\end{reponse}
\explication On compte les comparaisons dans le pire cas, un tableau trié à l'envers. Pour \code{i} = 1, la boucle \code{while} fait 1 tour ; pour \code{i} = 2, 2 tours ; \ldots ; pour \code{i} = $n-1$, $n-1$ tours. Au total : $1 + 2 + \cdots + (n-1) = \dfrac{n(n-1)}{2}$, de l'ordre de $n^2$. Vérification par programme : pour $n = 100$ éléments rangés à l'envers, on compte 4\,950 comparaisons, soit $100 \times 99 / 2$. (Dans le meilleur cas, un tableau déjà trié, il n'y a que $n - 1$ comparaisons ; la complexité retenue est celle du pire cas.)

\Qt{Question 3 -- Associer les fonctions aux courbes}{3 points}
\rappel{associer $10$, $x$, $\log_2(x)$, $x\log_2(x)$, $x^2$ et $2^x$ aux courbes A à F.}
\begin{reponse}
\begin{tabular}{@{}lll@{}}
A : $f(x) = 10$ & B : $f(x) = x$ & C : $f(x) = \log_2(x)$ \\
D : $f(x) = x\log_2(x)$ & E : $f(x) = x^2$ & F : $f(x) = 2^x$
\end{tabular}
\end{reponse}

\methode On calcule quelques valeurs simples et on les cherche sur le graphique.
\begin{center}\small
\begin{tabular}{|c|c|c|c|c|c|c|}\hline
$x$ & $10$ & $x$ & $\log_2(x)$ & $x\log_2(x)$ & $x^2$ & $2^x$ \\\hline
1 & 10 & 1 & 0 & 0 & 1 & 2 \\
2 & 10 & 2 & 1 & 2 & 4 & 4 \\
4 & 10 & 4 & 2 & 8 & 16 & 16 \\
5 & 10 & 5 & 2,32 & 11,6 & 25 & 32 \\
8 & 10 & 8 & 3 & 24 & 64 & 256 \\\hline
\end{tabular}
\end{center}
\begin{enumerate}
\item La \textbf{droite horizontale} est la fonction constante $10$.
\item La \textbf{droite qui passe par l'origine} et par (10~; 10) est $f(x) = x$.
\item La courbe qui monte \textbf{très lentement} (elle ne dépasse pas 5 sur tout le graphique) est $\log_2(x)$.
\item La courbe qui passe au-dessus de la droite $f(x)=x$ à partir de $x = 2$, et qui atteint 24 en $x = 8$, est $x\log_2(x)$.
\item Les deux courbes très raides se croisent en $x = 2$ et $x = 4$. Celle qui part de 1 en $x = 0$ et vaut 32 en $x = 5$ est $2^x$ ; celle qui part de 0 et vaut 25 en $x = 5$ est $x^2$.
\end{enumerate}

\begin{center}
\begin{tikzpicture}
\begin{axis}[width=0.9\linewidth,height=7cm,xmin=0,xmax=20,ymin=0,ymax=40,xlabel={$x$},ylabel={$f(x)$},grid=both,
 every axis plot/.append style={thick, black},domain=0.01:20,samples=200]
\addplot[solid] {10} node[pos=0.85, above, font=\small] {A : $10$};
\addplot[dashed] {x};
\addplot[dotted, very thick] {ln(x)/ln(2)} node[pos=0.85, above, font=\small] {C : $\log_2(x)$};
\addplot[dashdotted, domain=0.01:11.2] {x*ln(x)/ln(2)} node[pos=0.95, right, font=\small] {D : $x\log_2 x$};
\addplot[solid, domain=0:6.32] {x^2};
\addplot[densely dashed, domain=0:5.32] {2^x};
\draw[->, thin] (axis cs:2.6,35) node[left, font=\small, fill=white, inner sep=1pt] {F : $2^x$} -- (axis cs:5.0,32.5);
\draw[->, thin] (axis cs:8.2,27) node[right, font=\small, fill=white, inner sep=1pt] {E : $x^2$} -- (axis cs:5.5,29);
\node[font=\small, fill=white, inner sep=1pt] at (axis cs:17,22) {B : $x$};
\end{axis}
\end{tikzpicture}\\
{\small Les courbes du sujet, redessinées en noir et blanc, avec leur nom écrit à côté.}
\end{center}

\begin{remarque}
Dans le sujet, les courbes ne se distinguent que par leur couleur : imprimé en noir et blanc, le graphique ne permet pas de les reconnaître. De plus, les lettres A à F sont données dans le même ordre que la liste des fonctions, ce qui donne la réponse sans réfléchir. Ce corrigé explique comment reconnaître chaque courbe par ses valeurs.
\end{remarque}

\Qt{Question 4 -- Trace de \code{tri\_surprise([1, 5, 100, 78, 8, 22])}}{2 points}
\rappel{la valeur de \code{tab} à chaque itération de la boucle \code{while}.}

\methode Pour chaque \code{i}, on note \code{e = tab[i]}. Tant que l'élément à gauche (\code{tab[j]}) est plus grand que \code{e}, on le décale d'une case vers la droite. Puis on écrit \code{e} dans la case libre \code{tab[j+1]}.

\begin{reponse}
\begin{tabular}{|c|c|l|l|}\hline
\code{i} & \code{e} & Tour de \code{while} & \code{tab} \\\hline
1 & 5 & aucun ($5 > 1$) & \code{[1, 5, 100, 78, 8, 22]} \\\hline
2 & 100 & aucun ($100 > 5$) & \code{[1, 5, 100, 78, 8, 22]} \\\hline
3 & 78 & 1\textsuperscript{er} tour : 100 décalé & \code{[1, 5, 100, 100, 8, 22]} \\
  &    & fin : 78 inséré en case 2 & \code{[1, 5, 78, 100, 8, 22]} \\\hline
4 & 8 & 1\textsuperscript{er} tour : 100 décalé & \code{[1, 5, 78, 100, 100, 22]} \\
  &   & 2\textsuperscript{e} tour : 78 décalé & \code{[1, 5, 78, 78, 100, 22]} \\
  &   & fin : 8 inséré en case 2 & \code{[1, 5, 8, 78, 100, 22]} \\\hline
5 & 22 & 1\textsuperscript{er} tour : 100 décalé & \code{[1, 5, 8, 78, 100, 100]} \\
  &    & 2\textsuperscript{e} tour : 78 décalé & \code{[1, 5, 8, 78, 78, 100]} \\
  &    & fin : 22 inséré en case 3 & \code{[1, 5, 8, 22, 78, 100]} \\\hline
\end{tabular}\\[0.3em]
La fonction renvoie \code{[1, 5, 8, 22, 78, 100]}. La boucle \code{while} fait 5 tours en tout.
\end{reponse}
\explication Pendant les décalages, une valeur apparaît deux fois (par exemple \code{100, 100}) : c'est normal, la valeur de \code{e} est gardée dans la variable \code{e} et sera écrite à la fin. Pour \code{i} = 4, la boucle s'arrête quand \code{j} = 1, car \code{tab[1]} = 5 n'est pas plus grand que 8 : on écrit 8 dans \code{tab[2]}.

\begin{erreur}
Échanger les deux cases à chaque tour, comme dans un tri à bulles, ou oublier la dernière ligne \code{tab[j+1] = e}. Ici, on \textbf{décale} seulement, et \code{e} n'est écrit qu'une fois, à la fin.
\end{erreur}

\Qt{Question 5 -- Compléter la fonction \code{fusion}}{3 points}
\rappel{compléter les trous de \code{fusion(gauche, droite, tab)}.}
\begin{reponse}
\begin{lstlisting}
def fusion(gauche, droite, tab):
    i = 0
    j = 0
    while i + j < len(tab):
        if j == len(droite) or (i < len(gauche) and gauche[i] < droite[j]):
            tab[i + j] = gauche[i]
            i = i + 1
        else:
            tab[i + j] = droite[j]
            j = j + 1
\end{lstlisting}
\end{reponse}
\explication
\begin{itemize}
\item Lignes 2 et 3 : \code{i} est l'indice de lecture dans \code{gauche}, \code{j} dans \code{droite}. On commence au début des deux listes : 0.
\item Ligne 4 : \code{i + j} est le nombre d'éléments déjà recopiés, donc aussi la case de \code{tab} à remplir. On s'arrête quand \code{tab} est plein.
\item Ligne 5 : on prend dans \code{gauche} si \code{droite} est épuisée (\code{j == len(droite)}), ou s'il reste des éléments à gauche et que celui de gauche est le plus petit.
\item Lignes 6 et 7 : on recopie \code{gauche[i]} et on avance dans \code{gauche}. Lignes 9 et 10 : sinon, on recopie \code{droite[j]} et on avance dans \code{droite} (\code{j = j + 1}).
\end{itemize}

Trace de la dernière fusion de \code{tri\_fusion([1, 5, 100, 78, 8, 22])}, avec \code{gauche = [1, 5, 100]} et \code{droite = [8, 22, 78]} (déjà triées par les appels récursifs) :
\begin{center}\small
\begin{tabular}{|c|c|c|l|c|l|}\hline
Tour & \code{i} & \code{j} & Comparaison & Pris dans & \code{tab} en fin de tour \\\hline
1 & 0 & 0 & $1 < 8$ & gauche & \code{[1, \_, \_, \_, \_, \_]} \\
2 & 1 & 0 & $5 < 8$ & gauche & \code{[1, 5, \_, \_, \_, \_]} \\
3 & 2 & 0 & $100 > 8$ & droite & \code{[1, 5, 8, \_, \_, \_]} \\
4 & 2 & 1 & $100 > 22$ & droite & \code{[1, 5, 8, 22, \_, \_]} \\
5 & 2 & 2 & $100 > 78$ & droite & \code{[1, 5, 8, 22, 78, \_]} \\
6 & 2 & 3 & droite épuisée & gauche & \code{[1, 5, 8, 22, 78, 100]} \\\hline
\end{tabular}
\end{center}
(Les colonnes \code{i} et \code{j} donnent les valeurs au début du tour ; « \_ » marque une case pas encore remplie, qui contient encore l'ancienne valeur.) Le code complété a été testé sur 2\,000 tableaux tirés au hasard : il les trie tous correctement.

\begin{erreur}
Oublier d'avancer \code{j} (ou \code{i}) après avoir recopié un élément : \code{i + j} n'augmente plus et la boucle \code{while} ne s'arrête jamais. Autre erreur : commencer \code{i} ou \code{j} à 1, ce qui saute le premier élément de chaque liste.
\end{erreur}

\begin{remarque}
\code{tri\_fusion} trie \code{tab} \textbf{en place} (elle modifie la liste reçue). Elle ne renvoie \code{tab} que pour une liste de moins de 2 éléments ; sinon elle renvoie \code{None}. Ce n'est pas gênant ici, car les appels récursifs n'utilisent pas la valeur renvoyée, mais \code{print(tri\_fusion([3, 1, 2]))} affiche \code{None}.
\end{remarque}

\end{document}
